\chapter{Other Laboratories}
\label{ch:labs}

The queens are one hard counting problem among many, and the vertical view applies to all of them. This chapter runs through five others, the permanent, the Ising model, the volume of a convex body, satisfying assignments, and the number of linear extensions of a partial order, and for each describes the staircase, the level set sampler\index{level set sampler}, and the status of mixing. The pattern that emerges is uniform: in these five cases the problems for which the walk is known to be polynomial are exactly those whose decision versions are easy, and for every one of them the staircase has the same shape.

\section{The permanent}

For an $n \times n$ zero one matrix $A$ the permanent $\mathrm{per}(A) = \sum_\sigma \prod_i A_{i, \sigma(i)}$ counts the perfect matchings of the bipartite graph with adjacency matrix $A$. It is $\sharpP$ complete by Valiant\index{Valiant, Leslie}, and the existence of a perfect matching is in $\Pc$ by Edmonds\index{Edmonds, Jack}, and in $\RP$ by the determinant of the Edmonds matrix. The natural energy on permutations is $S(\sigma) = \#\{ i : A_{i, \sigma(i)} = 0 \}$, the number of non edges used, and the permanent is $N(0)$. The staircase has at most $n + 1$ steps.

Jerrum and Sinclair's 1989 algorithm runs a chain on matchings, not permutations, that adds or removes an edge or shifts an endpoint, and is reversible with respect to a distribution that weights matchings by a factor $\lambda^{\text{number of edges}}$; the ratio of perfect to near perfect matchings is estimated at a sequence of $\lambda$ values, which is an annealing\index{simulated annealing} ladder in $\lambda$. For dense matrices, with every row and column sum at least $n/2$, the chain mixes in polynomial time by canonical paths. Jerrum, Sinclair and Vigoda in 2004 handled arbitrary nonnegative matrices by weighting the near perfect matchings according to which vertices they miss, with the weights learned adaptively in a manner that resembles Wang Landau\index{Wang Landau algorithm}. The result is an FPRAS of running time $O^*(n^{10})$, later improved to $O^*(n^7)$ by Bez\'akov\'a, \v{S}tefankovi\v{c}, Vazirani and Vigoda using an improved cooling schedule.

The vertical reading: the ladder in $\lambda$ is a sequence of tilts of the staircase, and the adaptive weights of Jerrum, Sinclair and Vigoda are the flat histogram condition applied to the hole positions. The permanent\index{permanent} is the showcase for the thesis in the easy direction. Its decision problem is in $\Pc$, its staircase is short, its level sets are connected under local moves with polynomial congestion, and the walk is polynomial.

\section{The Ising model}

The ferromagnetic Ising model on a graph $G$ has configurations $\sigma \in \{\pm 1\}^V$, energy $H(\sigma) = -\sum_{(i,j) \in E} \sigma_i \sigma_j$, and partition function $Z(\beta) = \sum_\sigma e^{-\beta H(\sigma)}$. Computing $Z$ is $\sharpP$ complete even for planar graphs with a magnetic field, by Jerrum and Sinclair, and counting the ground states is the number of minimum cuts, which is in $\Pc$. The staircase has at most $|E| + 1$ steps indexed by the number of unsatisfied edges.

Jerrum and Sinclair gave an FPRAS\index{FPRAS} for $Z(\beta)$ at all $\beta$ in the ferromagnetic case, by running a chain not on spins but on the subgraphs of the high temperature expansion, which have no phase transition. Glauber dynamics mixes rapidly above the critical temperature and torpid\index{torpid mixing}ly below it on lattices, by the bottleneck between the two magnetised phases, while Swendsen\index{Swendsen Wang algorithm} Wang mixes rapidly for the Ising model on every graph at every temperature by Guo and Jerrum (2018); torpid mixing of Swendsen Wang is a Potts phenomenon, shown on the complete graph at criticality for $q \ge 3$ by Gore and Jerrum. The hard core model of independent sets is hard to approximate beyond the uniqueness threshold in the fugacity, by Sly's theorem of 2010, and the antiferromagnetic Ising model beyond its uniqueness threshold by Sly and Sun (2012); with Weitz's algorithm (2006) below the threshold this shows that the computational threshold coincides with the uniqueness threshold on the infinite regular tree.

The vertical reading: the staircase of the Ising model is smooth and well behaved, but its level sets are disconnected under spin flips below the critical temperature, which is the physical phase transition, and the walk needs the cluster or subgraph chains to get through. Sly's theorem is the sharpest known statement that a physical transition can be a computational one: on the hard side the staircase cannot be climbed, by any method, unless $\NP = \RP$.

\section{The volume of a convex body}

For a convex body $K \subseteq \mathbb{R}^n$ given by a membership oracle, the volume is $\sharpP$ hard to compute exactly and, by B\'ar\'any and F\"uredi, cannot be approximated by any deterministic polynomial time algorithm to within a factor better than $(cn / \log n)^{n/2}$. Dyer, Frieze and Kannan\index{Dyer, Frieze and Kannan} showed in 1991 that it can be approximated by a randomized algorithm, the first instance of a provable exponential gap between deterministic and randomized approximation. The algorithm is a ladder: a sequence of bodies $K_0 \subseteq K_1 \subseteq \dots \subseteq K_T = K$, with $K_0$ a ball of known volume and each $K_{i+1}$ at most a constant factor larger than $K_i$, and the ratios $\mathrm{vol}(K_i) / \mathrm{vol}(K_{i+1})$ estimated by sampling from $K_{i+1}$ with a random walk. The original running time $O^*(n^{23})$ has been reduced over thirty years to $O^*(n^4)$ by Lov\'asz\index{Lov\'asz, L\'aszl\'o} and Vempala and to $O^*(n^3)$ by Cousins and Vempala, through better walks, better isoperimetric inequalities, and better ladders, including the Gaussian cooling ladder of Cousins and Vempala in which $K_i$ is replaced by a Gauss\index{Gauss, Carl Friedrich}ian weighting of $K$ at inverse temperature $\beta_i$.

The vertical reading is exact. The likelihood is $L(x) = \Ind_K(x)$ under a Gaussian or uniform reference, the quantile function is a single step of width $\mathrm{vol}(K) / \mathrm{vol}(\text{reference})$, and the annealing ladder in $\beta$ is a sequence of tilts that walks the step from near one to its true position. The walk is polynomial because the level sets, which are convex bodies, are connected and have conductance bounded by the isoperimetric constant of a convex body, which Kannan, Lov\'asz and Simonovits showed is at least inverse polynomial. Applegate\index{Applegate, Kannan and Polson}, Kannan and Polson, and then Frieze, Kannan and Polson with the Metropolis\index{Metropolis algorithm} ball walk, extended the random walk and its analysis from indicators of convex bodies to log concave functions, which is the Bayesian case: a posterior proportional to a log concave likelihood times a log concave prior is sampled and integrated by the same walk, and the vertical ladder for its evidence has the same polynomial guarantee. This is a point at which Bayesian computation and complexity theory touched, and it is no accident that it was through a volume.

\section{Satisfying assignments}

For a Boolean formula $\varphi$ in conjunctive normal form, the number of satisfying assignments $\#\mathrm{SAT}$ is $\sharpP$ complete and the decision problem is $\NP$ complete. The natural energy is the number of unsatisfied clauses, and the staircase has at most $m + 1$ steps for $m$ clauses. A level set sampler would be a chain on assignments with at most $k$ unsatisfied clauses; WalkSAT and its relatives are such chains, used as optimisers rather than samplers, and their empirical success on random instances below the satisfiability threshold and failure above it is the heuristic form of a phase transition. The picture from the statistical physics of random constraint satisfaction, is that the solution space of a random $k$ SAT instance shatters into exponentially many clusters at a clause density below the satisfiability threshold, the clustering transition, and that local chains cannot move between clusters. This is the mechanism we conjectured for the queens in Chapter~\ref{ch:phase}, and in the random SAT setting with $k$ large enough it is established at the level of rigorous theorems by Achlioptas, Coja Oghlan and others.

The vertical reading: the staircase of $\#\mathrm{SAT}$ is short and the walk needs $O(m)$ rungs, but the bottom rungs are level sets that shatter into exponentially many components, and no local chain mixes across them. The clustering transition is the mixing conjecture made concrete, and it is a theorem for random instances with $k$ large enough.

\section{Linear extensions}

A partial order on $n$ elements has some number of linear extensions, the total orders compatible with it, and Brightwell and Winkler showed in 1991 that counting them is $\sharpP$ complete. Finding one is trivial by topological sort. The problem is self reducible\index{self reducibility} and the level sets are connected: the linear extensions index the simplices of volume $1/n!$ in the canonical triangulation of the order polytope, and a chain that swaps adjacent elements when the swap respects the order mixes in polynomial time by a coupling\index{coupling} argument due to Bubley and Dyer, with mixing time $O(n^3 \log n)$. The count therefore has an FPRAS, and indeed the volume of the order polytope, which is the count divided by $n!$, can also be approximated by the convex body algorithm.

The vertical reading: the energy is the number of order violations, the ground step is the set of linear extensions, and the walk is polynomial because the adjacent transposition chain on the ground step itself mixes rapidly. This is a case where the bottom rung, which for the queens is the hardest, is directly samplable, and the ladder is not needed; the count can be obtained by the Jerrum Valiant Vazirani self reducibility route from a uniform sampler on the ground step alone.

\section{The pattern}

Set the five problems and the queens side by side.

The permanent, the ferromagnetic Ising model, the volume and linear extensions have decision versions in $\Pc$, staircases of polynomial length, and level sets whose connectivity and conductance\index{conductance} can be bounded (for the Ising model by the cluster or subgraph chains and not by spin flips), and the walk is polynomial for all of them. Satisfying assignments and the queens completion have $\NP$ complete decision versions, staircases of polynomial length, and level sets that are believed, and for random $k$ SAT with $k$ large proved, to shatter at low energy, and the walk is believed to be exponential for both.

The staircase itself is never the obstruction. In every case it has polynomially many steps, and in every case the quantile reordering reduces the exponential sum to a polynomial sum of widths. The obstruction, when there is one, is always the bottom rungs, and always for the same reason: the level sets fall apart into pieces that no efficiently implementable chain can connect. The thesis of this book is that this pattern is not a coincidence but the content of $\NP \neq \RP$ seen from the quantile axis. The same axis, with a single step, is where $\Pc = \RP$ lives, and the question there is whether the single bottom rung, which has density at least a half, can be hit without a chain at all.
